\subsection{Conventions and the actual manuscript claims}

We use the contravariant crystalline Dieudonn\'e functor $\mathbb D$.
Thus its Hodge sequence is
\[
 0\longrightarrow\omega_G\longrightarrow\mathbb D(G)_S
 \longrightarrow\operatorname{Lie}(G^D)\longrightarrow0.
\]
Here $G^D$ is the Cartier dual and $\omega_G=e^*\Omega^1_{G/S}$.
All finite group schemes below are commutative; the Cartier dual and the
Dieudonn\'e functors used in the question require this hypothesis.
We use cohomological degrees, so the co-Lie complex
$\ell_G=Le^*L_{G/S}$ has amplitude $[-1,0]$ for a finite locally free group.
We visually checked the original archive images of cote 129, scans
94, 114, and 115, in addition to the supplied transcription.
Scan 115 states V \emph{up to isogeny}. Scan 94 writes the finite module
$M=D^*(G)$ in its derived tensor formula; it does not put $\Delta^*(G)$
in that formula. These distinctions affect the answers.

\subsection{B.1: the exact integral and rational answer}

For a $p$-adically separated complete ring $A$, put $A_0=A/pA$.
Evaluation on the compatible canonical PD thickenings
$A/p^nA\twoheadrightarrow A_0$ defines the finite projective module
\[
 P_A(G_0)=\varprojlim_n\mathbb D(G_0)_{A/p^nA}.
\]
The integral target category has objects $(G_0,L)$, where $G_0$ is a
Barsotti--Tate group over $A_0$ and $L\subset P_A(G_0)$ is a direct summand
lifting its Hodge filtration. A morphism $(G_0,L_G)\to(H_0,L_H)$ is a
homomorphism $f_0:G_0\to H_0$ for which
$\mathbb D(f_0)_A(L_H)\subseteq L_G$. In the rational target we replace
group homomorphisms by $\operatorname{Hom}[1/p]$, the evaluated modules by
$P_A[1/p]$, and the filtrations by their scalar extensions.
These conventions specify precisely the two full-faithfulness assertions.

\begin{lemma}[Proved: canonical divided powers, including torsion]
\label{B:PD-bounds}
Every $p$-adically complete ring $A$ has the canonical divided powers
\[
 \gamma_m(pa)=\frac{p^m}{m!}a^m\quad(m\geq1)
\]
on $pA$. They descend to $A/p^NA$. If $p\geq3$, these divided powers
are nilpotent modulo $p^N$: every divided-power monomial of total weight
at least $2N$ is zero. For $p=2$, their restriction to $4A$ is nilpotent
modulo $2^N$, with weight $N$ sufficient. The divided powers on
$2\mathbb Z_2/2^N\mathbb Z_2$ need not be nilpotent when $N\geq2$.
\end{lemma}

\begin{proof}
All coefficients $p^m/m!$ lie in $p\mathbb Z_p$. If $pa=pb$, put
$c=a-b$. Since $pc=0$ and $a^m-b^m$ is a multiple of $c$, the displayed
formula has the same value for $a$ and $b$. The divided-power identities
are polynomial identities over $\mathbb Z_p$ in these coefficients;
they hold in the torsion-free universal polynomial ring and therefore
after specialization to $A$. For $x=p^Na$, its $m$th divided power is
$p^{Nm}a^m/m!$, which belongs to $p^NA$, so the quotient is defined.

For $p\geq3$,
$v_p(p^m/m!)=m-v_p(m!)\geq m/2$. A product of divided powers of
elements of $pA$, with indices summing to $r$, therefore belongs to
$p^{\lceil r/2\rceil}A$. For $4A$ we have
$v_2(4^m/m!)=2m-v_2(m!)=m+s_2(m)\geq m$, where $s_2(m)$ is the sum
of the binary digits. The same multiplication argument proves the
second bound. Finally,
$v_2(2^{2^r}/(2^r)!)=1$ for all $r$. These divided powers of the single
element $2$ remain nonzero modulo $4$, so no nilpotence bound exists.
\end{proof}

\begin{lemma}[Known, with the needed uniform bound proved]
\label{B:square-zero-Hom}
Let $B\twoheadrightarrow R$ have kernel $I$ with $I^2=pI=0$, and assume
that $p$ is nilpotent in $B$. For Barsotti--Tate groups $G,H$ over $B$,
\[
 \operatorname{Hom}_B(G,H)\longrightarrow
 \operatorname{Hom}_R(G_R,H_R)
\]
is injective and its cokernel is killed by $p$.
\end{lemma}

\begin{proof}
We use two precisely stated standard facts: for every $B$-algebra $C$,
$H(C)\to H(C/IC)$ is surjective, and its kernel is killed by $p$.
These are \cite[Lemmas 3.2 and 3.3]{B:LauRelations}; the latter uses the
formal smoothness of a Barsotti--Tate group as a sheaf. Their hypotheses
hold after this base change because $(IC)^2=pIC=0$.
For $f_0:G_R\to H_R$ and $x\in G(C)$, choose a lift
$y\in H(C)$ of $f_0(\bar x)$ and define $f(x)=py$.
Different lifts differ by a point killed by $p$, so this is independent
of the choice. It is additive and natural in $C$, hence defines a
homomorphism of fppf sheaves, and therefore of Barsotti--Tate groups.
Its reduction is $pf_0$.
If a homomorphism $h:G\to H$ reduces to zero, the same kernel assertion
gives $ph=0$. Since $[p]:G\to G$ is an epimorphism of fppf sheaves and
$h\circ[p]=ph$, we obtain $h=0$.
This proves both claims, with an exponent independent of $G,H,C$.
\end{proof}

\begin{lemma}[Known, with the nonnoetherian comparison explained]
\label{B:effectivity}
For a $p$-adically separated complete ring $A$,
\[
 \operatorname{BT}(A)\longrightarrow
 \varprojlim_N\operatorname{BT}(A/p^NA)
\]
is an equivalence. No noetherian or bounded-$p$-torsion hypothesis is
required in this affine assertion.
\end{lemma}

\begin{proof}
The proof of \cite[Lemma 2.4.4]{B:deJong1995} constructs the group one
finite flat level at a time. We recall why it applies here. Finite
projective $A$-modules are equivalent to compatible systems of finite
projective $A/p^N A$-modules: lift an idempotent matrix for the first
module successively across nilpotent ideals, take its limit in a matrix
ring over $A$, and identify the given lifts using the usual lifting of
isomorphisms of projectives. Morphisms are inverse limits because
finite projectives are direct summands of finite free modules.
For each $r$, take the inverse limit of the finite projective Hopf
algebras of $G_N[p^r]$. The multiplication, comultiplication, identity,
and antipode are morphisms between finite projective modules, and hence
lift uniquely. Their identities hold modulo every $p^N$, and hence
over $A$ by separatedness. The transition maps lift in the same way.

The finite-flat exactness conditions defining a Barsotti--Tate system
also transfer. For a transition morphism, write $B$ for the lifted
finite projective coordinate algebra of the target. It is $p$-adically
complete. The coordinate algebras of the source at every level are
compatible finite projective $B/p^N B$-modules. The preceding
idempotent lifting argument, now applied over $B$, proves that their
limit is finite projective over $B$. Its positive local ranks modulo
$p$ imply faithful flatness, since $pB$ lies in the Jacobson radical.
Equality with the prescribed
kernel is an isomorphism of finite presented affine schemes modulo
every $p^N$, and the associated maps of finite projective coordinate
modules are isomorphisms by Nakayama. Thus the lifted levels form a
Barsotti--Tate group. The argument for homomorphisms is the same inverse
limit argument for Hopf algebra maps. This proves essential
surjectivity and full faithfulness.
\end{proof}

\begin{theorem}[Proved: Nice V and the complete integral criterion]
\label{B:NiceV}
Let $A$ be any $p$-adically separated complete commutative ring. Then
\[
 G\longmapsto(G_0,\omega_G\subseteq P_A(G_0))
\]
is fully faithful up to isogeny for every prime $p$. Integrally, on the
whole category $\operatorname{BT}(A)$, it is fully faithful if and only
if either $p\geq3$, or $p=2$ and $2A=0$.
For odd $p$ it is in fact an equivalence onto the category of pairs
defined above. Consequently R3's odd-prime assertion and its reduction
at $2$ are correct, and the first assertion of L5 is correct.
\end{theorem}

\begin{proof}[Proof of the positive assertions]
For odd $p$, Lemma~\ref{B:PD-bounds} allows us to apply the
Grothendieck--Messing deformation theorem to
$A/p^NA\twoheadrightarrow A_0$ for every $N$. We use its equivalence
of categories, including morphisms, not merely a bijection on
isomorphism classes; the exact form is recalled in
\cite[Remark 2.6]{B:LauSmoothness}, with references to Messing V.1.6
and the comparison with the BBM crystal. Evaluating the crystal and
the direct summand modulo $p^N$ gives exactly its deformation data.
The uniqueness in that equivalence makes the deformations compatible
as $N$ varies. Lemma~\ref{B:effectivity} then gives the asserted
equivalence over $A$.

For $p=2$, let $f_0:G_0\to H_0$ be an integral homomorphism whose
evaluated crystalline map respects the integral Hodge summands.
Lemma~\ref{B:square-zero-Hom}, applied to $A/4A\to A/2A$, gives a
homomorphism $f_1:G_{A/4A}\to H_{A/4A}$ reducing to $2f_0$.
For $N\geq2$ apply Grothendieck--Messing to
$A/2^NA\to A/4A$ with its divided powers on $4A/2^NA$.
The base-change isomorphism of the Dieudonn\'e crystal identifies
its evaluation here with $P_A(G_0)/2^N$ and identifies
$\mathbb D(f_1)$ with $2\mathbb D(f_0)$. This comparison comes from the
map of PD thickenings induced by the inclusion $4A\subset2A$; both
PD structures are restrictions of the same canonical one.
Thus $f_1$ preserves the Hodge filtration and has a unique lift at
every level. Compatibility and Lemma~\ref{B:effectivity} give a
homomorphism over $A$ lifting $2f_0$.

For a rational target morphism, first clear the denominator of $f_0$.
The failure to preserve the integral Hodge summands is a map between
finite projective $A$-modules which becomes zero after inverting $p$.
Its matrix has finitely many entries, so one common power of $p$
annihilates it. Multiplying $f_0$ by that power puts us in the preceding
integral situation. At $2$ one additional multiplication by $2$
suffices. Dividing the lifted morphism by the total power gives the
required rational morphism.

Reduction of homomorphisms is injective: at each finite level factor
$A/p^N\to A/p$ into the square-zero steps
$A/p^j\to A/p^{j-1}$ and apply
Lemma~\ref{B:square-zero-Hom}; then use separatedness and
Lemma~\ref{B:effectivity}. This proves faithfulness, including after
inverting $p$. If $2A=0$, reduction is the identity on the group
itself, so integral full faithfulness is tautological.
\end{proof}

\begin{proposition}[Proved: the characteristic-two obstruction for every $A$]
\label{B:Kummer}
If $A$ is $2$-adically separated complete and $2A\ne0$, the functor in
Theorem~\ref{B:NiceV} is not full integrally. The counterexample is the
base change of two ordinary height-two Barsotti--Tate groups over
$\mathbb Z_2$.
\end{proposition}

\begin{proof}
For a unit $q\in\mathbb Z_2^\times$ and $m=2^n$, define
\[
 E_q[m]=\coprod_{i=0}^{m-1}
 \operatorname{Spec}\bigl(\mathbb Z_2[X]/(X^m-q^i)\bigr).
\]
Write a point as $(x,i)$. Its group law is
\[
 (x,i)(y,j)=
 \bigl(xyq^{-\lfloor(i+j)/m\rfloor},\ i+j\bmod m\bigr).
\]
The carry identity for addition modulo $m$ proves associativity;
the identity is $(1,0)$ and the inverse for $i>0$ is $(q/x,m-i)$.
The coordinate algebras are monic of degree $m$, hence finite free,
and there is an fppf exact sequence
$0\to\mu_m\to E_q[m]\to\mathbb Z/m\to0$.
The embeddings $(x,i)\mapsto(x,2i)$ give a Barsotti--Tate system.
Indeed multiplication by $2$ from level $2m$ to level $m$ is
\[
 (x,i)\longmapsto
 \bigl(x^2q^{-\lfloor i/m\rfloor},\ i\bmod m\bigr);
\]
it is fppf surjective by extraction of a square root, and its kernel
is $E_q[2]$. These formulas also identify every $E_q[m]$ with the
$m$-torsion of the resulting sheaf $E_q$.

Take $q=1$ and $q=-1$. Both reductions modulo $2$ identify canonically
with $\mu_{2^\infty}\oplus\mathbb Q_2/\mathbb Z_2$.
If $q^a=r^b$, the formulas
\[
 (x,i)\longmapsto
 \bigl(x^a r^{-\lfloor bi/m\rfloor},\ bi\bmod m\bigr)
\]
give compatible homomorphisms $E_q\to E_r$, acting by $a$ on the
multiplicative part and by $b$ on the \`etale quotient.
In particular $u:E_1\to E_{-1}$ with $(a,b)=(1,2)$ and
$v:E_{-1}\to E_1$ with $(a,b)=(2,1)$ have composites $[2]$.

The Hodge lines of these two lifts coincide under their common
crystalline identification. To see this integrally, work first over
$\mathbb Z_2$. In the splitting of the special-fibre crystal into
multiplicative and \`etale lines, $\mathbb D(u)$ has matrix
$\operatorname{diag}(1,2)$. The induced map on invariant
differentials is an isomorphism, since $u$ is the identity on the
formal multiplicative subgroup. The image of the Hodge line of
$E_{-1}$ therefore equals the multiplicative line, which is the
Hodge line of $E_1$. A generator $(a,b)$ of the former satisfies
$2b=0$ and $a\in\mathbb Z_2^\times$. Thus $b=0$ and the two lines
are equal. Equality survives arbitrary base change to $A$.

Their common identity modulo $2$ is consequently a morphism in the
target category. If it lifted to $E_{1,A}\to E_{-1,A}$, the level-two
point $(1,1)$ would map to $(r,1)$ with $r^2=-1$ and $r\equiv1\pmod{2A}$.
The component really is $1$: idempotents do not change across an
ideal in the Jacobson radical, and $2A$ lies in that radical.
Write $r=1+2b$. Then
\[
 0=r^2+1=2(1+2b+2b^2).
\]
The second factor is a unit, again because $2A$ is in the Jacobson
radical. Hence $2=0$, a contradiction. This argument also applies to
rings containing all $2$-power roots of $-1$: the congruence required
for lifting this particular identity is essential.

The finite obstruction modulo $4$ and all carry identities through
$m=64$ were checked by \texttt{code/B\_finite\_certificates.py}; the
general identities and the finite flat/BT hypotheses were proved
above. This completes the necessity in Theorem~\ref{B:NiceV}.
\end{proof}

\begin{proposition}[Known classification, with the DVR deduction made explicit]
\label{B:DVR-classification}
Let $A=\mathcal O_K$ be a complete discrete valuation ring of mixed
characteristic $(0,p)$ with perfect residue field $k$. The integral
full-faithfulness assertion of Theorem~\ref{B:NiceV} for odd $p$, and
its rational assertion for every $p$, also follow directly from the
Breuil--Kisin classification. At $p=2$ integral fullness fails for
every such $A$ by Proposition~\ref{B:Kummer}.
\end{proposition}
\begin{proof}
Put $W=W(k)$, choose a uniformizer $\pi$ and its Eisenstein polynomial
$E(u)\in W[u]$, and let $S$ be the $p$-adic completion of the
divided-power envelope of $(E(u))\subset W[u]$. Its quotient by
$\operatorname{Fil}^1 S$ is $A$. Give $S$ the Frobenius extending
Witt Frobenius and sending $u$ to $u^p$. It is $p$-torsion free.
For a Barsotti--Tate group $G/A$, the classifying object is
\begin{align*}
 M_G&=\mathbb D(G_0)(S),\\
 \operatorname{Fil}^1M_G&=
  \{m\in M_G:m\bmod\operatorname{Fil}^1S\in\omega_G\},\\
 \varphi_1&=\varphi/p\quad\hbox{on }\operatorname{Fil}^1M_G.
\end{align*}
These are exactly Kisin's Appendix, Lemma (A.2) and definition (A.5).
His Proposition (A.6), pp.~31--33 of the opened author preprint,
states that this contravariant functor is an equivalence for $p>2$
and an equivalence on isogeny categories for $p=2$
\cite{B:KisinCrystals}. Breuil's earlier integral equivalence for odd
$p$ is \cite[Theorem 4.2.2.9]{B:Breuil2000}.

An integral $f_0:G_0\to H_0$ preserving the Hodge summands gives,
by crystalline evaluation, a Frobenius-compatible $S$-linear map
$M_H\to M_G$. The map $S\to A$, $u\mapsto\pi$, identifies its
reduction with the evaluated map in the question, so the displayed
preimage description shows that it preserves
$\operatorname{Fil}^1$. Because both modules are $p$-torsion free,
commutation with $\varphi$ implies commutation with
$\varphi_1=\varphi/p$. Thus it is a morphism of the classifying
objects. The same reasoning after inverting $p$ applies to a rational
target morphism. Proposition (A.6) gives the required morphism of
groups, integrally for odd $p$ and rationally for all $p$.

Its reduction really is the prescribed $f_0$. Their difference has
zero crystalline evaluation on $S$, hence also on $W$ via $u\mapsto0$.
Classical Dieudonn\'e theory over the perfect field $k$ makes its
restriction to $k$ zero. The map
$\operatorname{Hom}(G_0,H_0)\to\operatorname{Hom}(G_k,H_k)$
is injective by Lemma~\ref{B:square-zero-Hom}, applied successively
to the nilpotent ideal $(\pi)$ of $A/p$; this remains injective
after inverting $p$. Consequently the reductions agree. The same
argument proves faithfulness of the construction.

Kisin assumes a perfect residue field, explicitly on the first page
of the preprint. Theorem~\ref{B:NiceV}, whose proof uses no residue
field assumption, supplies the answer for complete discrete
valuation rings with imperfect residue field as well. We do not
silently apply the perfect-field classification to those rings.
\end{proof}

\subsection{B.2: non-nilpotent divided powers}

\begin{theorem}[Known: both alternatives in Nice IV are valid]
\label{B:NiceIV}
Let $B\twoheadrightarrow R=B/I$ be a PD extension of rings in which
$p$ is nilpotent. No nilpotence condition on the divided powers is
assumed. The functor sending a Barsotti--Tate group over $B$ to its
reduction over $R$ and its lifted Hodge direct summand is an
equivalence on each of the following categories:
\begin{enumerate}
 \item groups with connected geometric fibres;
 \item groups whose geometric fibres have no nonzero multiplicative
 part, equivalently groups whose Cartier duals have connected fibres.
\end{enumerate}
If the divided powers are nilpotent, the equivalence holds without
either group-theoretic condition. These assertions hold at $p=2$
and do not require noetherian bases. The corresponding scheme
statements follow by Zariski descent.
\end{theorem}

\begin{proof}
This is exactly \cite[Remark 2.6]{B:LauSmoothness}, whose statement
and proof we opened. Its ``infinitesimal'' groups are the groups
with connected fibres in (1), and its ``unipotent'' groups are those
in (2). The source uses the covariant crystal; applying it to the
Cartier dual and then taking the orthogonal Hodge summand gives the
contravariant formulation here. Cartier duality exchanges the two
conditions, and the crystalline duality comparison used in that
remark is BBM, Section 5.3.

For clarity, the reason that arbitrary divided powers are allowed
under these conditions is part of the cited proof. For the formal
group relevant to a unipotent Barsotti--Tate group, the augmentation
ideal $J$ of its Cartier-dual algebra is a nilideal. Divided powers
on $I$ induce nilpotent divided powers on $IJ$. The formal logarithm
and exponential in the deformation proof are therefore inverse
even though the divided powers on $I$ itself are not nilpotent.
The source then invokes the full deformation theorem and verifies
that the comparison with the BBM crystal extends. Thus this is an
equivalence with the same Hodge-filtration target as in the question,
not merely an existence statement for a display lift.
Its \cite[Corollary 2.7]{B:LauSmoothness} identifies this crystal with
the crystal of the nilpotent display when the PD structure is
compatible with the canonical divided powers of $p$; compare also
\cite[Corollary 97]{B:ZinkDisplay}. The latter compatibility is needed
for that display comparison, not for the more general Remark 2.6
just used. Hence the second assertion of L5 is correct with the
precise fibrewise meanings stated above.
\end{proof}

\subsection{B.3: ordinary crystals, admissibility, and counterexamples}

Write $\operatorname{DC}(S)$ for ordinary finite locally free crystals
on $(S/\mathbb Z_p)_{\mathrm{cris}}$, with operators
$F:\varphi^*\mathcal M\to\mathcal M$ and
$V:\mathcal M\to\varphi^*\mathcal M$ satisfying $FV=VF=p$.
No divided Frobenius is included in this definition.
A filtered Dieudonn\'e crystal has in addition a locally direct
summand $L\subset\mathcal M_S$ with the following admissibility
condition. For every characteristic-$p$ PD object $(U,T,\delta)$,
Frobenius on $T$ factors as $\varphi_{U/T}:T\to U$, and
\begin{equation}\label{B:admissibility}
 \varphi_{U/T}^*(L|_U)=
 \ker\bigl(F:\varphi^*\mathcal M_{(U,T)}\to\mathcal M_{(U,T)}\bigr).
\end{equation}
This is \cite[Section 3, equation (3.1)]{B:LauDivided}.
It specifies the essential-image condition; merely imposing
$FV=VF=p$ is not a substitute for it on singular bases.

\begin{theorem}[Known: verified positive classes and essential images]
\label{B:positive-crystals}
The ordinary crystalline Dieudonn\'e functor is fully faithful in
each of the following cases:
\begin{enumerate}
 \item $S$ is locally the spectrum of a quasisyntomic
 $\mathbb F_p$-algebra;
 \item $S$ locally admits a $p$-basis, without a finiteness restriction
 on that basis or a noetherian hypothesis;
 \item $S$ is an excellent locally noetherian local complete
 intersection scheme;
 \item $S$ is reduced, $F$-finite, $F$-nilpotent, and PD torsion free.
\end{enumerate}
Here $F$-finite means that absolute Frobenius is finite;
$F$-nilpotent means that its kernel ideal is locally nilpotent;
PD torsion free means that the $p$-complete universal PD envelopes
of affine presentations are $p$-torsion free.
In (1), the essential image consists exactly of the crystals admitting
the filtration \eqref{B:admissibility}; this filtration is unique and
every crystal morphism preserves it. In (2) every object of
$\operatorname{DC}(S)$ belongs to the essential image.
In (4) the essential image is again exactly the crystals admitting
\eqref{B:admissibility}. The same essential-image assertion holds for
$F$-finite local complete intersections.
\end{theorem}

\begin{proof}
For (1), \cite[Theorem 4.74]{B:ALB} gives the antiequivalence with
admissible prismatic Dieudonn\'e crystals. We spell out the transfer
to the target in the question. By its Proposition 3.21 one can choose
a quasisyntomic cover with all \v Cech terms quasiregular semiperfect.
In characteristic $p$ the initial prism is the crystalline PD ring
$A_{\mathrm{cris}}$, and it is $p$-torsion free
(Section 3.4 and Lemma 3.27). Thus on every such term a crystalline
Dieudonn\'e module is the same module with Frobenius and
$V=pF^{-1}$; the comparison for the group functors is Theorem 4.44.
Proposition 4.29 identifies admissible objects fully faithfully with
windows over the Nygaard frame. Remark 4.30 identifies those windows
with precisely the admissible Hodge filtrations
\eqref{B:admissibility}. Descent for the modules, filtrations, and
groups therefore identifies the classical target with this full
subcategory. This proves full faithfulness, uniqueness of the
filtration, and the asserted whole essential image. This is also
the conclusion explicitly recorded in its Remark 1.19; the
admissibility qualification must be retained.

For (2) use \cite[Theorem 7.6]{B:LauDivided}; its Lemma 3.3 says
that every ordinary crystal has a unique admissible filtration.
For (3) use \cite[Theorem 4.6]{B:deJongMessing1999}, whose
local-complete-intersection and excellence hypotheses are exactly
those stated here. This theorem asserts full faithfulness; it is
not being used to assert essential surjectivity beyond its statement.
For (4), \cite[Theorem 7.9]{B:LauDivided} gives an equivalence with
filtered Dieudonn\'e crystals, and its Lemma 3.2 makes forgetting
the filtration fully faithful on a reduced scheme. The same lemma
and Corollary 7.10 prove the assertion for $F$-finite local complete
intersections. These are direct compositions of the cited functors,
so both their variance and their essential images agree with the
claims above.
\end{proof}

\begin{remark}[Known: what the stronger modern target changes]
\label{B:divided-target}
Lau's \cite[Theorem 9.3]{B:LauDivided} is an equivalence on every
$F$-finite, $F$-nilpotent characteristic-$p$ scheme, in particular
every $F$-finite noetherian scheme, with \emph{divided} Dieudonn\'e
crystals. The latter retain a divided Frobenius $\Phi_1$ on the
filtration. When the PD rings have $p$-torsion, $p\Phi_1=\Phi$ does
not determine $\Phi_1$. The theorem is therefore not a theorem about
the ordinary target of scan 114. Its Corollary 9.5 does give the
ordinary functor full faithfulness up to isogeny; more precisely,
the integral Hom map is injective with cokernel killed by $p$.
The counterexamples below exhibit why this qualification matters.
\end{remark}

\begin{lemma}[Proved: a finite PD detector, including the prime $2$]
\label{B:finite-PD}
For every prime $p$ put
\[
 B_p=\mathbb F_p[x,y,u,v]/
 (x^3,x^2y,xy^2,y^3,xu,yu,xv,yv,u^2,v^2)
\]
and $K=(x^2,xy,y^2,u,v)$. There is a divided-power structure on $K$
with $\gamma_p(x^2)=u$, $\gamma_p(y^2)=v$, and
$\gamma_p(xy)=0$. In particular the image of
$\gamma_p(x^2)\gamma_p(y^2)-\gamma_p(xy)^2$ is the nonzero element $uv$.
\end{lemma}

\begin{proof}
The displayed quotient has basis
$1,x,y,x^2,xy,y^2,u,v,uv$; in particular $uv\ne0$.
For odd $p$ define, on
$z=ax^2+bxy+cy^2+du+ev+fuv$,
\[
 \delta_p(z)=a^pu+c^pv.
\]
For $p=2$ define instead
\[
 \delta_2(z)=a^2u+c^2v+deuv.
\]
These maps satisfy the divided-$p$ conditions of
\cite[Lemma 23.5.3]{B:StacksPD}. Indeed $K^3=0$, $xK=yK=0$, and
$K^2$ is spanned by $uv$. For odd $p$ the additive cross terms of
total degree $p$ vanish, and $z^p=0$. At $2$ the additive identity
is $\delta_2(z+w)=\delta_2(z)+\delta_2(w)+zw$; its only nonzero
cross term is $(de'+d'e)uv$, supplied by the last term in the
formula. Scalar compatibility follows because multiplication by
$a_0+a_1x+a_2y+k$ on $K$ changes its $x^2,xy,y^2,u,v$
coefficients by $a_0$ and changes only its $uv$ coefficient further.
The square of a scalar at $2$ differs from $a_0^2$ by a combination
of $x^2,y^2$, which annihilates the image of $\delta_2$.
For odd $p$ its $p$th power is $a_0^p$.
The cited lemma therefore supplies the full divided powers.

At $2$ one can avoid even this existence criterion: take
$\gamma_1(z)=z$, $\gamma_2(z)=\delta_2(z)$,
$\gamma_3(z)=z\delta_2(z)$, $\gamma_4(z)=\delta_2(\delta_2(z))$,
and $\gamma_n(z)=0$ for $n\geq5$. The preceding multiplication rules
verify the identities directly. The attached program exhaustively
verifies them on all $512$ elements of $B_2$ and all $64$ elements
of $K$, including every nonzero index range. It checks $32768$
additivity identities, $131072$ scalar identities, and $1024$ each
of product and composition identities.

The formula for $\delta_p$ displayed in
\cite[Example 60.22.1]{B:StacksCounterexamples} omits the term $deuv$.
That formula is correct for odd $p$, but at $2$ it fails already
on $u+v$: the divided-square identity requires
$\gamma_2(u+v)=uv$. The correction above is needed to validate
the characteristic-two instance of that source's counterexample.
\end{proof}

\begin{proposition}[Refuted: the unrestricted conjecture of scan 114]
\label{B:crystal-counterexamples}
For every prime $p$, the ordinary crystalline Dieudonn\'e functor
fails to be full over the following bases:
\begin{align*}
 R_0&=\mathbb F_p[x,y]/(x^2,xy,y^2),\\
 R&=\mathbb F_p[a,b,c,d]/(ac,ad,bc,bd),\\
 R_\infty&=R[t_1,t_2,t_3,\ldots].
\end{align*}
The first base is finite and nonreduced; the second is reduced and
of finite type; the third is reduced and nonnoetherian.
In all three cases a nonzero morphism
\[
 \mathbb D(\mu_{p^\infty})\longrightarrow
 \mathbb D(\mathbb Q_p/\mathbb Z_p)
\]
does not come from a homomorphism
$\mathbb Q_p/\mathbb Z_p\to\mu_{p^\infty}$.
\end{proposition}

\begin{proof}
For $R_0$ take the $p$-complete PD envelope $D_0$ of its displayed
presentation over $\mathbb Z_p$, relative to the canonical PD on
$(p)$, and put
\[
 \tau_0=\gamma_p(x^2)\gamma_p(y^2)-\gamma_p(xy)^2.
\]
For $R$ take the analogous envelope $D$ and put
\[
 \tau=\gamma_p(ac)\gamma_p(bd)-\gamma_p(ad)\gamma_p(bc).
\]
We verify all needed properties, rather than infer them from the
existence of torsion somewhere in a PD ring. First, multiplication
by $p!$ makes either difference zero, by
$\gamma_p(rs)=r^p\gamma_p(s)$ and
$(ac)(bd)=(ad)(bc)$. Since $(p-1)!$ is a unit, $p\tau=0$,
and similarly $p\tau_0=0$.
Second, the PD derivation of $\tau$ is zero. For example its
$da$ coefficient is
\[
 \frac{a^{p-1}}{(p-1)!}
 \bigl(c^p\gamma_p(bd)-d^p\gamma_p(bc)\bigr)=0,
\]
since the two terms in parentheses are both $\gamma_p(bcd)$.
The $db,dc,dd$ coefficients vanish by the same displayed identity
with the letters interchanged. Substitution $a=c=x$, $b=d=y$
then proves $d\tau_0=0$ as well.
The crystalline PD de Rham comparison
\cite[Proposition 60.21.3 and Example 60.22.1]{B:StacksCounterexamples}
therefore makes these elements global sections of the crystalline
structure sheaf, equivalently scalar endomorphisms of that crystal.

Frobenius is represented on the envelopes by raising the polynomial
variables to their $p$th powers. For every PD element $z$,
\[
 \gamma_p(z^p)=(p!)^{p-1}\gamma_p(z)^p.
\]
Hence $\varphi(\tau)$ is
$(p!)^{2p-2}$ times the difference of the $p$th powers of the two
products defining $\tau$. That difference is a multiple of $\tau$.
The coefficient is divisible by $p$, so $\varphi(\tau)=0$.
The same proof applies to $\tau_0$.
Lemma~\ref{B:finite-PD} gives a PD map $D_0\to B_p$ carrying
$\tau_0$ to $uv\ne0$. Thus $\tau_0\ne0$. The substitution just
used extends by the universal property to $D\to D_0$ and carries
$\tau$ to $\tau_0$, proving $\tau\ne0$ too.

For the contravariant convention,
$\mathbb D(\mu_{p^\infty})=\mathcal O_{\mathrm{cris}}$ has
$F=p$, $V=1$, whereas
$\mathbb D(\mathbb Q_p/\mathbb Z_p)=\mathcal O_{\mathrm{cris}}$
has $F=1$, $V=p$. Multiplication by either scalar is therefore a
morphism of Dieudonn\'e crystals: both compatibility equations
are exactly $p\tau=\varphi(\tau)$. It is nonzero by the detector.

There is no corresponding group homomorphism. On a reduced
characteristic-$p$ ring every $p$-power root of unity is $1$.
On $R_0$, such a root lies in $1+(x,y)$, and this group is killed
by $p$. A homomorphism from the constant group
$\mathbb Q_p/\mathbb Z_p$ would give roots $\zeta_n$ satisfying
$\zeta_n=\zeta_{n+1}^p$, so all are $1$ in this case too.
This proves the asserted zero Hom groups.

The ideal defining $R$ is $(a,b)\cap(c,d)$, an intersection of
prime ideals, so $R$ is reduced. It describes two affine planes
meeting at a point. Polynomial extension preserves reducedness,
and the ascending chain $(t_1)\subset(t_1,t_2)\subset\cdots$
proves that $R_\infty$ is nonnoetherian. Pull back the nonzero
crystalline scalar from $R$ to $R_\infty$. Setting all $t_i=0$
is a retraction, so its pullback cannot be zero. All operator
identities persist under base change, and the group Hom is still
zero by reducedness. This proves the third example.

The finite $p=2$ detector, including the correction essential to
that prime, is exhaustively certified in
\texttt{logs/B\_finite\_certificates.jsonl}. The proofs above work
for every prime; the computation is not claimed to enumerate an
infinite family of primes.
\end{proof}

\begin{remark}[Proved consequence: reducedness and finiteness are not criteria]
\label{B:base-kinds}
Each of the three kinds requested in B.3 contains both positive
and negative examples. A field is a positive reduced example;
$\mathbb F_p[\epsilon]/\epsilon^2$ is a positive nonreduced example
by Theorem~\ref{B:positive-crystals}(1); and the perfect ring
$\mathbb F_p[t^{1/p^\infty}]$ is a positive nonnoetherian example
by (2). The preceding proposition supplies the negative examples.
Consequently none of ``reduced'', ``nonreduced'', or
``nonnoetherian'' can replace the hypotheses in the positive
theorems. We searched for the reduced two-plane construction in
the original crystalline literature and in searches combining
``Dieudonn\'e'', ``reduced'', ``not fully faithful'', and
``two planes''. We found the Berthelot--Ogus/Stacks scalar example
and the general torsion obstruction, but no exact match for the
displayed reduced reduction. We claim no priority.
\end{remark}

\subsection{B.4: the complex on every base and the corrected tensor formula}

\begin{theorem}[Proved: construction of the manuscript's complex]
\label{B:Delta}
For a commutative finite locally free group scheme $G$ over a scheme
$S$, there is a contravariant perfect complex $\Delta_S(G)$ of Tor
amplitude $[0,1]$, compatible with arbitrary base change, and a
functorial distinguished triangle
\begin{equation}\label{B:Delta-triangle}
 \ell_G[-1]\longrightarrow\Delta_S(G)
 \longrightarrow\ell_{G^D}^{\vee}\longrightarrow\ell_G.
\end{equation}
In particular it exists on every base on which $p$ is locally
nilpotent, as asked on cote 129, scan 94.
\end{theorem}

\begin{proof}
Raynaud's embedding theorem gives, Zariski locally on $S$, an exact
sequence of fppf sheaves represented by group schemes
\[
 0\longrightarrow G\longrightarrow A\xrightarrow{q}B\longrightarrow0
\]
with $A,B$ abelian schemes and $q$ an isogeny. We use the exact
statement in \cite[Theorem 4.65]{B:ALB}, which cites BBM, Theorem
3.1.1. Define on such an open set
\begin{equation}\label{B:Delta-resolution}
 C(A,q)=
 [H^1_{\mathrm{dR}}(B/S)\xrightarrow{q^*}
  H^1_{\mathrm{dR}}(A/S)]
 \quad\text{in degrees }0,1.
\end{equation}
Both terms are vector bundles commuting with arbitrary base change,
and their Hodge sequences are functorial exact sequences
$0\to\omega_A\to H^1_{\mathrm{dR}}(A/S)\to
\operatorname{Lie}(A^D)\to0$.
These standard abelian-scheme facts are recalled, with the BBM
reference, in \cite[Proposition 4.57]{B:ALB}.
Consequently there is a termwise exact sequence of complexes
\[
 0\to[\omega_B\to\omega_A]\to C(A,q)
 \to[\operatorname{Lie}(B^D)\to\operatorname{Lie}(A^D)]\to0,
\]
all written in degrees $0,1$.

The square defining $G$ as the inverse image of the identity of $B$
is Tor independent because $q$ is finite flat. Cotangent-complex
base change and transitivity therefore identify
$\ell_G$ with $[\omega_B\to\omega_A]$ in degrees $-1,0$.
Dualizing the isogeny gives
$0\to G^D\to B^D\to A^D\to0$; the same argument, followed by
duality of perfect complexes, identifies the rightmost complex
with $\ell_{G^D}^{\vee}$ (with the usual sign in the dual
differential). The short exact sequence yields
\eqref{B:Delta-triangle} locally.

We give the descent argument, since merely choosing an abelian
embedding would not define a global object. A morphism between
two such presentations which induces the identity on $G$ gives,
by pullback of de Rham cohomology, a map between their complexes
in the reverse direction. The preceding cotangent identifications
are natural for this map and identify both outer terms of the
triangle with the identity. Hence the map of the middle complexes
is a quasi-isomorphism. Two presentations have a common refinement:
embed $G$ diagonally in $A_1\times A_2$ and take its abelian
quotient. Its projections give morphisms to both presentations.
Products of any finite list give these refinements associatively,
and the projection maps strictly commute.

Here is an explicit way to retain all descent coherence. Choose
an open cover $U_i$ with presentations. On an intersection
$U_{i_0}\cap\cdots\cap U_{i_r}$ use the product presentation
$G\hookrightarrow\prod_j A_{i_j}$ and its complex. For an
inclusion of lists use the pullback map induced by the projection
of products. These maps are quasi-isomorphisms by the preceding
paragraph, and their compositions agree, so they form a derived
\v Cech descent datum. Totalizing its \v Cech complex, with derived
extension from the intersections, produces a complex on $S$.
Its restriction to $U_i$ is quasi-isomorphic to $C(A_i,q_i)$:
adjoining the index $i$ gives the usual extra-index contraction
of the augmented \v Cech diagram, and every comparison involved
is a quasi-isomorphism. Thus this totalization is locally a
two-term vector-bundle complex, proving perfection and the stated
Tor amplitude. Applying the identical construction to the short
exact sequence of complexes glues the distinguished triangle.
Refining a cover or adding another set of presentations gives
the same comparison by products and the same contraction;
therefore it does not change the resulting object or triangle.

For completeness, a homomorphism $f:G\to H$ also gives the
contravariant map asserted in the theorem. If $G\hookrightarrow A$
and $H\hookrightarrow C$ are presentations, replace the first
embedding by its graph $G\hookrightarrow A\times C$ using $f$.
Projection to $C$ extends $f$ and induces a map of quotients, hence
a chain map $C(C,C/H)\to C(A\times C,(A\times C)/G)$.
Projection to $A$ identifies the latter complex with the one for
$G$. This defines $\Delta(H)\to\Delta(G)$.
Two extensions of $f$ between fixed presentations differ by a map
$A\to C$ vanishing on $G$; it factors uniquely through $B=A/G$.
Pullback by that factor is the chain homotopy between the two
maps of two-term de Rham complexes. Products compare arbitrary
choices, and products of three presentations prove compatibility
with composition. The same products used in the \v Cech
construction give its higher coherence. All formulas commute
with arbitrary base change, since abelian quotients, the
cotangent complexes in question, and $H^1_{\mathrm{dR}}$ do.
This completes the construction and all its comparisons.
\end{proof}

\begin{proposition}[Proved: the perfect-field formula with its normalization]
\label{B:coLie-Dieudonne}
Let $k$ be a perfect field of characteristic $p$, put $W=W(k)$, and
let $\sigma$ be Witt-vector Frobenius. Write
\[
 \mathcal D=W\langle F,V\rangle/
 (FV=VF=p,\ Fa=\sigma(a)F,\ Va=\sigma^{-1}(a)V).
\]
For a finite commutative group scheme of $p$-power order, let $M$
be its classical contravariant left Dieudonn\'e module, normalized
as on scan 94. If $P$ denotes its crystalline Dieudonn\'e module,
then $P=\sigma^*M$. With this convention there are natural
isomorphisms
\begin{align}
 \ell_G&\simeq(\mathcal D/F\mathcal D)
                 \otimes^{\mathbf L}_{\mathcal D}M,
                 \label{B:coLie-formula}\\
 \Delta_k(G)&\simeq
       \bigl(\sigma^*M\otimes_W^{\mathbf L}k\bigr)[-1].
                 \label{B:Delta-field}
\end{align}
The quotient in the first line is a \emph{right} $\mathcal D$-module.
\end{proposition}

\begin{proof}
The normalization $P=\sigma^*M$ is the classical-to-crystalline
comparison: see \cite[Proposition 4.51]{B:ALB} for Barsotti--Tate
groups and \cite[Theorem 3.8]{B:MondalClassifying} for finite groups
(the latter explicitly states BBM, Theorem 4.2.14, with this twist).
This is why a twist must not be silently dropped for a general
perfect field.

First take a Barsotti--Tate group $H$ over $k$.
For its finite free crystalline module $P_H$, the usual Hodge
condition is $\sigma^*\omega_H=\ker(F\bmod p)$ inside
$\sigma^*P_H/p$. Since $FV=VF=p$ and $P_H$ is free over the DVR
$W$, the matrix factorization gives
$\operatorname{im}(V\bmod p)=\ker(F\bmod p)$ and
$\ker(V\bmod p)=\operatorname{im}(F\bmod p)$.
Thus $P_H/FP_H\simeq\sigma^*\omega_H$, or
$M_H/FM_H\simeq\omega_H$ in the classical normalization.

Left multiplication by $F$ on $\mathcal D$ is injective. Indeed
the usual normal form embeds this ring into the skew Laurent ring
$W[1/p][F,F^{-1};\sigma]$ by $V=pF^{-1}$.
Consequently
$[\mathcal D\xrightarrow{F\cdot}\mathcal D]$ in degrees $-1,0$
is a free right-module resolution of $\mathcal D/F\mathcal D$.
Its left $W$-module structure in degree $-1$ is twisted so that
the displayed map is $W$-linear; after tensoring it is the
linearization $\sigma^*M\xrightarrow{F}M$.
For $M=M_H$ this map is injective, again because $VF=p$ and
$M_H$ is $p$-torsion free. Its cokernel is $\omega_H$.

Choose an abelian presentation $0\to G\to A\to B\to0$.
Classical Dieudonn\'e theory gives the exact sequence
\[
 0\to M_{B[p^\infty]}\to M_{A[p^\infty]}\to M_G\to0.
\]
Applying the derived right-module tensor functor just computed
shows that its value on $M_G$ is
$[\omega_B\to\omega_A]$ in degrees $-1,0$, which is $\ell_G$
by the cotangent calculation in Theorem~\ref{B:Delta}.
The comparison is independent of the presentation because all
maps in that construction and in this derived triangle are natural.
This proves \eqref{B:coLie-formula}.
For the crystalline modules the corresponding two free $W$-modules
resolve $P_G=\sigma^*M_G$. Derived reduction modulo $p$ and a shift
by $[-1]$ therefore give precisely the two-term de Rham complex
\eqref{B:Delta-resolution}. This proves \eqref{B:Delta-field}.
\end{proof}

\begin{proposition}[Refuted: the tensor assertion of L6]
\label{B:L6-refutation}
The existence and triangle assertion in L6 are true, but its formula
with $\Delta(G)$ in place of $M$ is false, already for
$G=\alpha_2$ over $\mathbb F_2$. The corrected statement is
Theorem~\ref{B:Delta} over a general base and
Proposition~\ref{B:coLie-Dieudonne} over a perfect field.
\end{proposition}

\begin{proof}
At $k=\mathbb F_2$ there is no Frobenius-twist ambiguity.
For $G=\alpha_2$, $M=k$ with $F=V=0$. Formula
\eqref{B:coLie-formula} gives $[k\xrightarrow0 k]$ in degrees
$-1,0$, also obtained directly from the equation $T^2=0$ at its
identity. Formula \eqref{B:Delta-field} gives
$[k\xrightarrow0 k]$ in degrees $0,1$.
Apply $(\mathcal D/F\mathcal D)\otimes^{\mathbf L}_{\mathcal D}-$
to this latter complex, using the right-module resolution from
the preceding proof. Its differentials are all zero, and its
cohomology dimensions in degrees $-1,0,1$ are respectively
$1,2,1$. They differ from those of $\ell_{\alpha_2}$, which are
$1,1,0$. Thus even the natural perfect-field interpretation of
the proposed tensor expression fails. Over a general base the
expression additionally needs a specified ring sheaf and module
actions; the field counterexample already refutes its assertion.

The order-$p$ tests, in the classical normalization, are
\[
\begin{array}{c|c|c|c}
G&M&(F,V)&\ell_G\text{ in degrees }[-1,0]\\\hline
\mu_p&k&(0,\sigma^{-1})&[k\xrightarrow0 k]\\
\mathbb Z/p&k&(\sigma,0)&[k\xrightarrow{\sim}k]\simeq0\\
\alpha_p&k&(0,0)&[k\xrightarrow0 k].
\end{array}
\]
In each case $\Delta$ has one-dimensional cohomology in degrees
$0$ and $1$; the maps in its triangle distinguish the cases.
The finite complex calculation for $\alpha_2$ is included in
\texttt{code/B\_finite\_certificates.py} and its output log.
\end{proof}
